% 第 3 章 · API 直收组 —— Copilot Chat（微软/GitHub）
\section{Copilot Chat（微软/GitHub）v0.48.1 —— API 直收：服务端点菜}
\label{sec:copilot}

\begin{verdictbox}{vcoral}
确认隐藏上传：external ingest 管线把变更文件全文 POST 到 \ep{api.github.com/external/code/ingest*}；\ep{/batch} 由服务端自选 \ep{doc\_ids} 点菜回传；门槛是 onExp 隐藏设置，未设值用户可经 ExP-TAS 中会话热翻转；consent modal 双实现零调用点。
\end{verdictbox}

微软自家一方管线，也是九家中披露面最厚的——但披露厚度与门槛诚实度不匹配。与 ZCode 打包整树走对象存储不同，这条链不产本地归档：客户端上交全库文件指纹，服务端 \ep{/batch} 逐批回传想要的 \ep{doc\_ids}，客户端只对被点名的读盘发全文。

点菜权连同门本身都在服务端：设置 \ep{chat.workspace.codeSearchExternalIngest.enabled} 声明 \ep{default:false}、\ep{tags:[advanced,onExp]}（extension/package.json:1），用户未显式设值时取值回落到 ExP-TAS 实验变量（treatment variable）——微软可在零用户操作下远程翻转这道门。

\begin{figure}[!t]
\centering
\includegraphics[width=\linewidth,height=0.82\textheight,keepaspectratio]{f02-copilot.pdf}
\caption{Copilot Chat v0.48.1 external ingest 泳道。控制面 ExP-TAS 30 min 轮询可远程翻转隐藏门；服务端 \ep{/ingest/batch} 自选 \ep{doc\_ids} 点菜，\ep{/document} 逐文件发全文；consent modal 双实现零调用者（虚线幽灵框）。}
\label{fig:copilot-pipe}
\end{figure}

\subsection{机制}

门控两段式（图~\ref{fig:copilot-pipe}）。\fn{isExternalIngestEnabled}（extension.js:2387）读 \fn{getExperimentBasedConfig}（:7115）：用户显式值优先，未设值依次回落 \ep{experimentName}、\ep{copilotchat.config.<id>}、\ep{config.<fullyQualifiedId>} 三个 treatment variable，最后 default。

experimentName 绑定 \ep{copilotchat.config.chat.advanced.workspace.codeSearchExternalIngest.enabled}（:1195）；TAS 客户端每 30 min 拉 \ep{default.exp-tas.com/vscode/ab} 的 \ep{Configs[].Parameters}（:1113），遥测关闭则整条通道停用。取 \ep{"force"} 时 ingest 取代 remote-index 路径；该门独立于可见的 \ep{chat.workspace.enableCodeSearch}——external ingest 单独开启即可令语义搜索可用（isAvailable）。

触发无定时器，仅两个调用点：每次 @workspace 语义搜索进 \fn{searchLocalDiff}（:2409），先 \fn{updateForceIncludeFiles} 再 \fn{doIngest}，progress 传 noop；diff 超 2000 文件或逾 70\% 标 tooLarge 时跳过前者、后者照跑。另一路是 Developer 命令 \ep{github.copilot.buildRemoteWorkspaceIndex}（:6211）。

收集与上传全在 \fn{ExternalIngestIndex}/\fn{ExternalIngestClient}（:2363--2386）：\fn{reconcileDbFiles} 对每 folder 跑 \ep{findFiles(**/*)} 入 sqlite \ep{codebase-external.sqlite}，并注册 FileSystemWatcher 持续跟踪。

\fn{performIngestion} 四步 POST——\ep{/external/code/ingest} 交 \ep{fileset\_name=vscode.copilot-chat.<uuid>}、\ep{new\_checkpoint}、\ep{geo\_filter}（wasm 概率集）、\ep{coded\_symbols} 换 \ep{ingest\_id}；\ep{/coded\_symbols} 区间循环对账；\ep{/batch} 分页取服务端自选 \ep{doc\_ids}；每 \ep{doc\_id} 读盘 base64 POST \ep{/document}，64 并发、500 重试 3 次、429 重试 10 次，读失败仍发空文档；\ep{/finalize} 封口。

凭证为 \fn{getGitHubSession} 静默取回的用户 OAuth Bearer（\ep{silent:!0}），附 \ep{X-Client-Application} 等三头——明文 JSON over TLS，无 STS、无客户端加密、无独立对象存储段。

\subsection{载荷}

服务端分三段实收：握手段拿到 fileset uuid、全树 checkpoint、覆盖全部 docSha 的 GeoFilter 与符号清单，等价于完整文件清单指纹；\ep{/document} 段收到被点名文件的 b64 全文与 \ep{file\_path=<folderId>/<相对路径>}（folderId 消毒后截 8 字符）；查询段 \ep{/external/embeddings/code/search} 把 \ep{prompt}（$\le$7800B 原文）绑 \ep{scoping\_query=fileset:<uuid>}，模型 \ep{metis-1024-I16-Binary}。

量级：5000 文件首全量约 55MB；增量只发被点名 \ep{doc\_ids}（20 文件约 220KB）；checkpoint 零漂移 0 字节。

卫生面是四层 JS 名单加 wasm：扩展名黑名单 177 条含全套私钥材料（\ep{pfx}/\ep{p12}/\ep{pem}/\ep{crt}/\ep{cer}/\ep{key}/\ep{priv}/\ep{jks}/\ep{keystore}/\ep{csr}），目录段 26 条、scheme 11 条，basename 仅 5 条（\ep{.ds\_store}、\ep{thumbs.db}、\ep{package-lock.json}、\ep{yarn.lock}、\ep{.cache}）——\ep{.env}、\ep{id\_rsa}、\ep{.npmrc}、\ep{credentials.json}、\ep{netrc}、\ep{.ssh} 全数放行；wasm IngestFilter 14 条 deny 正则同样不含密钥基名（extension.js:1854、wasm bindings :265）。

\ep{.gitignore} 只靠 VS Code 默认 \ep{search.useIgnoreFiles} 兜底；Copilot 内容排除 \fn{isCopilotIgnored} 生效（服务端 \ep{/copilot\_internal/content\_exclusion} 下发）。

\subsection{consent 面}

用户可见面三处：状态项 \ep{copilot.workspaceIndexStatus} 写 "Indexes your codebase for more relevant AI results"（:6214）、build/delete 两条命令、设置描述句 "Enable external ingest for semantic codebase search."——管线全链用户可见文案无一处 "upload"。

唯一的 consent 对话 \fn{promptForExpandedLocalIndexing} 有两处实现（:2363 stub \ep{return!1}、:6982 真 modal "Build local index for this workspace?"），在 20MB bundle 内调用点数为 0。auth 全程静默，上传进度 reporter 是 noop。

宣称对照见表~\ref{tab:copilot-claims}。唯一能对上的是 repository-indexing 文档页 "This feature uploads your data to GitHub to make it searchable"——它把范围框在 non-GitHub repositories；同页 "controlled by policy and is disabled by default" 对个人用户不成立，门槛是服务端可翻的实验变量而非组织策略。

产品内 aka.ms 直达页写 "parts of the index might come from remote sources"，数据流向说反。GitHub 自家 Extension Developer Policy 明令禁未经 consent 收集个人数据、禁误导——第三方扩展按同一标尺会被下架。

\begin{table}[!t]
\centering\small
\begin{tabularx}{\linewidth}{@{}p{0.40\linewidth}Xc@{}}
\toprule
宣称原文（逐字） & 实际行为 & 判定 \\
\midrule
"This feature uploads your data to GitHub to make it searchable." & 与 \ep{/external/code/ingest*} 逐字吻合；仅框在 non-GitHub repositories 语境 & 如实 \\
"controlled by policy and is disabled by default" & 个人用户门槛 = onExp 设置 + TAS treatment variable，可远程热翻 & 低估 \\
"parts of the index might come from remote sources" & 方向说反：本地文件内容上传到 GitHub & 误导 \\
"Indexes your codebase for more relevant AI results."（状态栏） & 管线全链文案无 "upload"；consent modal 死代码 & 误导 \\
\bottomrule
\end{tabularx}
\caption{宣称对照摘录（disclosure/gh-repository-indexing.md、vsc-workspace-context.md 逐字）。}
\label{tab:copilot-claims}
\end{table}

\subsection{时间线}

v0.8.0（2023-10-05）至 v0.45.1（2026-04-23）九个二分探针的 5 组特征字面量全部零命中；v0.45.1 是最后一个可获取无签名版本。marketplace 1125 版列表在 v0.45.1 到 v0.48.1 之间零条目，v0.45.2/v0.46.0/v0.47.0/v0.48.0 及 3 个日期预发布号共 7 探针全 404——首发真身被 404 窗口遮蔽，v0.48.1（2026-05-15）只是首个可获取携带版，且为可获取历史最新版（v0.48.2 起探针亦 404）。

落地即完全体：onExp 设置、TAS 翻转接线、端点族、sqlite、consent modal 死代码同版到位。开源对照：vscode-copilot-chat 仓的门是 \ep{ConfigKey.TeamInternal}，其 package.json 无此设置项——shipped 版把内部开关外化成隐藏 contributed setting。

\evid{evidence/copilot-chat/figure-pack.json · notes/sweep/copilot-chat.md · \ep{extracted/copilot-chat/extension/dist/extension.js:2363/2409/7115/1854/6214/6982}}
